% !TeX spellcheck = en_GB
\ifcsname SlidesDistr\endcsname%
\documentclass[handout,aspectratio=169]{beamer}
\else%
\documentclass[aspectratio=169]{beamer}
\fi%
\input{../syde556_lecture_slides_preamble}

\date{October 27, 2026}
\title{SYDE 556/750 \\ Simulating Neurobiological Systems \\ Lecture 7: Temporal Basis Functions}

\begin{document}
	
	\begin{frame}{}
		\vspace{0.5cm}
		\begin{columns}[c]
			\column{0.6\textwidth}
			\MakeTitle
			\column{0.4\textwidth}
			\includegraphics[width=\textwidth]{media/a22_2007_06_einsteinturm_sonnenuhr_p6070062.jpg}
		\end{columns}
	\end{frame}

	\begin{frame}{Representing Stimulus Histories}
		\centering%
		\includegraphics[width=\textwidth]{media/delay_network_decoding.pdf}%
	\end{frame}

	\begin{frame}{Representing Functions: Sampling}
		\centering%
		\includegraphics<1>[width=\textwidth]{media/fun_1_q4_1.pdf}%
		\includegraphics<2>[width=\textwidth]{media/fun_1_q4_2.pdf}%
		\includegraphics<3>[width=\textwidth]{media/fun_1_q4_3.pdf}%
		\includegraphics<4>[width=\textwidth]{media/fun_1_q4_4.pdf}%
		\includegraphics<5>[width=\textwidth]{media/fun_1_q4_5.pdf}%
		\includegraphics<6>[width=\textwidth]{media/fun_1_q4_6.pdf}%
%		\includegraphics<7>[width=\textwidth]{media/fun_1_q6_1.pdf}%
		\includegraphics<7>[width=\textwidth]{media/fun_1_q6_2.pdf}%
		\includegraphics<8>[width=\textwidth]{media/fun_1_q6_3.pdf}%
		\includegraphics<9>[width=\textwidth]{media/fun_1_q6_4.pdf}%
		\includegraphics<10>[width=\textwidth]{media/fun_1_q6_5.pdf}%
		\includegraphics<11>[width=\textwidth]{media/fun_1_q6_6.pdf}%
  \end{frame}
  \begin{frame}{Nyquist Sampling}
		\centering%
    \begin{mdframed}
      \emph{The Nyquist-Shannon Sampling Theorem}\\
      If $f(t)$ contains no frequencies greater than $B$ then it is \emph{completely} determined by samples spaced $\Delta t = \frac{1}{2B}$ apart ($N = 2BT$ equally spaced samples for a time-slice $[0, T)$). There is a \emph{one-to-one mapping} between the samples $\vec x$ and the function $f(t)$.
    \end{mdframed}
	\end{frame}
  \begin{frame}{Representing Functions: Sampling <3Hz}
		\centering%
		\includegraphics<1>[width=\textwidth]{media/fun_2_q6_1.pdf}%
		\includegraphics<2>[width=\textwidth]{media/fun_2_q6_2.pdf}%
		\includegraphics<3>[width=\textwidth]{media/fun_2_q6_3.pdf}%
		\includegraphics<4>[width=\textwidth]{media/fun_2_q6_4.pdf}%
		\includegraphics<5>[width=\textwidth]{media/fun_2_q6_5.pdf}%
		\includegraphics<6>[width=\textwidth]{media/fun_2_q6_6.pdf}%
		\includegraphics<7>[width=\textwidth]{media/fun_2_q6_small.pdf}%
	\end{frame}

	\begin{frame}{Representing Functions: Fourier Basis}
		\includegraphics<1>[width=\textwidth]{media/fourier_q1.pdf}%
		\includegraphics<2>[width=\textwidth]{media/fourier_q2.pdf}%
		\includegraphics<3>[width=\textwidth]{media/fourier_q3.pdf}%
		\includegraphics<4>[width=\textwidth]{media/fourier_q4.pdf}%
		\includegraphics<5>[width=\textwidth]{media/fourier_q5.pdf}%
		\includegraphics<6>[width=\textwidth]{media/fourier_q6.pdf}%
	\end{frame}

	\begin{frame}{Representing Functions: Cosine Basis}
		\includegraphics<1>[width=\textwidth]{media/cosine_q1.pdf}%
		\includegraphics<2>[width=\textwidth]{media/cosine_q2.pdf}%
		\includegraphics<3>[width=\textwidth]{media/cosine_q3.pdf}%
		\includegraphics<4>[width=\textwidth]{media/cosine_q4.pdf}%
		\includegraphics<5>[width=\textwidth]{media/cosine_q5.pdf}%
		\includegraphics<6>[width=\textwidth]{media/cosine_q6.pdf}%
	\end{frame}

	\begin{frame}{Representing Functions: ReLU Basis}
		\includegraphics<1>[width=\textwidth]{media/neural_q1.pdf}%
		\includegraphics<2>[width=\textwidth]{media/neural_q2.pdf}%
		\includegraphics<3>[width=\textwidth]{media/neural_q3.pdf}%
		\includegraphics<4>[width=\textwidth]{media/neural_q4.pdf}%
		\includegraphics<5>[width=\textwidth]{media/neural_q5.pdf}%
		\includegraphics<6>[width=\textwidth]{media/neural_q6.pdf}%
	\end{frame}

	\begin{frame}{Representing Functions: Legendre Basis}
		\includegraphics<1>[width=\textwidth]{media/legendre_q1.pdf}%
		\includegraphics<2>[width=\textwidth]{media/legendre_q2.pdf}%
		\includegraphics<3>[width=\textwidth]{media/legendre_q3.pdf}%
		\includegraphics<4>[width=\textwidth]{media/legendre_q4.pdf}%
		\includegraphics<5>[width=\textwidth]{media/legendre_q5.pdf}%
		\includegraphics<6>[width=\textwidth]{media/legendre_q6.pdf}%
	\end{frame}

	\begin{frame}{Implementing the Delay Network}
		\centering%
		\begin{columns}
			\column{0.5\textwidth}
			\centering
			\includegraphics{media/lti_integrator_vs_neural_a.pdf}%
			\column{0.5\textwidth}
			\centering
			\includegraphics{media/lti_integrator_vs_neural_b.pdf}%
		\end{columns}
		\begin{align*}
			\theta \mat{A} &= a_{ij} \in \mathbb{R}^{q \times q} \,, & \kern-0.5em a_{ij} &= \begin{cases}
			(2i + 1)(-1) & \kern-0.5em i < j \,,\\
			(2i + 1)(-1)^{i - j + 1} & \kern-0.5em i \geq j
			\end{cases}\\
			\theta \mat{B} &= b_i \in \mathbb{R}^{q} \,, & \kern-0.5em b_i &= (2i + 1)(-1)^i \\
			\mat A' &= \tau \mat A + \mat I & \mat B' &= \tau \mat B
		\end{align*}
	\end{frame}

	\begin{frame}{Delay Network: Step Function}
		\centering%
		\includegraphics<1>[width=\textwidth]{media/delay_network_step_q1.pdf}%
		\includegraphics<2>[width=\textwidth]{media/delay_network_step_q2.pdf}%
		\includegraphics<3>[width=\textwidth]{media/delay_network_step_q3.pdf}%
		\includegraphics<4>[width=\textwidth]{media/delay_network_step_q4.pdf}%
		\includegraphics<5>[width=\textwidth]{media/delay_network_step_q5.pdf}%
		\includegraphics<6>[width=\textwidth]{media/delay_network_step_q6.pdf}%
		\includegraphics<7>[width=\textwidth]{media/delay_network_step_q11.pdf}%
	\end{frame}

	\begin{frame}{Delay Network: Windowed Sine Function}
		\centering%
		\includegraphics<1>[width=\textwidth]{media/delay_network_sine_i0_q6.pdf}%
		\includegraphics<2>[width=\textwidth]{media/delay_network_sine_i1_q6.pdf}%
		\includegraphics<3>[width=\textwidth]{media/delay_network_sine_i2_q6.pdf}%
	\end{frame}

	\begin{frame}{From the Delay Network to the Legendre Memory Unit}
		\begin{columns}
			\column{0.55\textwidth}
			\begin{itemize}
				\item<1-> The \emph{Legendre Memory Unit} (LMU) places the delay network's linear system at the core of a recurrent neural architecture
				\item<2-> Parameter-efficient: $\mathcal{O}(N)$ rather than the expected $\mathcal{O}(N^2)$
				\item<3-> Trainable feedforward, deployable recurrently
				\item<4-> Neural connection: explains time cells; runs dynamic and spiking (or not)
			\end{itemize}
			\column{0.45\textwidth}
			\centering
			\includegraphics[width=\textwidth]{media/lmu_architecture.png}
		\end{columns}
	\end{frame}

	\begin{frame}{State Space Models (SSMs)}
		\begin{itemize}
			\item<1-> LSTMs (1995): long the dominant recurrent network, but training cannot be parallelized, and sequences beyond ${\sim}10^3$ samples are out of reach
			\item<2-> Transformers (2017): parallel training, but compute and memory scale as $\mathcal{O}(N^2)$ in sequence length
			\item<3-> \emph{SSMs}: the delay network/LMU (2018--2019, Waterloo) was the first; re-derived at Stanford as HiPPO (2020), followed by the S-series (2021) and Mamba (2023)
			\item<4-> SSMs now achieve state-of-the-art results on language, audio, and genomics, with SSM work at DeepMind, Meta, Nvidia, Princeton, Stanford, \dots
		\end{itemize}
	\end{frame}

	\begin{frame}{LMU Benchmarks}
		\begin{columns}
			\column{0.55\textwidth}
			\centering
			\includegraphics[width=0.82\textwidth]{media/lmu_memory_capacity.png}
			\column{0.4\textwidth}
			\centering
			\includegraphics[height=0.52\textheight]{media/lmu_psmnist_table.png}
		\end{columns}
		\begin{itemize}
			\item Memory: ${\sim}10^6\times$ lower error while remembering ${\sim}10^4\times$ longer sequences than an LSTM
			\item psMNIST: (simplified) LMU set the state of the art (98.49\,\%) with 60\,\% fewer parameters
			\item Scaling: $\mathcal{O}(N)$ compute and $\mathcal{O}(1)$ memory, vs.\ $\mathcal{O}(N^2)$ for transformers
		\end{itemize}
	\end{frame}

	\begin{frame}{SSMs in Practice (and Other Bases)}
		\begin{columns}
			\column{0.47\textwidth}
			\begin{itemize}
				\item<1-> Applications: LLMs (e.g.\ Mamba), speech recognition, keyword spotting, RF classification, anomaly detection, biosignals, nonlinear system prediction
				\item<2-> Other temporal bases work too (Fourier, cosine, Haar, \dots); analytic bases often beat learned ones
			\end{itemize}
			\column{0.53\textwidth}
			\centering
			\includegraphics<2->[width=\textwidth]{media/lmu_bases_boxplot.png}
		\end{columns}
		\onslide<3->
		\begin{mdframed}
			``Next gen state space models [are] probably the next generation of transformers.''\hfill --- Jensen Huang, CEO Nvidia, 2024
		\end{mdframed}
	\end{frame}

	\backupbegin

	\begin{frame}[noframenumbering]{Image sources}
		\small
		\textbf{Title slide}\\Infrared Photograph of a Sundial Near the Einstein Tower in Potsdam, Germany\\Author: DrNRNowaczyk, 2007.\\From \href{https://commons.wikimedia.org/wiki/File:A22_2007-06_Einsteinturm_Sonnenuhr_p6070062.jpg}{Wikimedia}.
	\end{frame}


	\backupend
	
\end{document}
